\section{Background}
\label{sec:background}

\subsection{The Radar Receive Chain}
Fig.~\ref{fig:processing-chain} shows an exemplar receive chain considered in this review. After digital downconversion,
the data for one coherent processing interval (CPI) form a complex-valued cube of $N_s$ fast-time samples, $N_p$ pulses,
and $N_c$ receive channels, each element an in-phase/quadrature (IQ) sample \cite{richardsPreviewBasicRadar2022}.
Each stage then operates along one or more axes of this cube. Pulse compression applies a matched filter along fast
time, normally as a fast convolution. Doppler processing applies an FFT along slow time; in a row-major layout this
needs a corner turn (transpose) to keep memory access contiguous. Beamforming combines the channel axis, either with
fixed weights or adaptively through Space-Time Adaptive Processing (STAP), which estimates a covariance matrix and solves a linear system for each range cell. Finally, CFAR detection compares each cell with a noise estimate drawn from 
its neighbouring cells in a sliding window formation to set an adaptive detection threshold. In SAR, image formation 
uses FFT-based algorithms or backprojects every pulse onto every pixel to generate an image \cite{richardsPreviewBasicRadar2022}.

\begin{figure}[!t]
    \centering
    \resizebox{\columnwidth}{!}{\begin{tikzpicture}[
    node distance = 6mm and 5mm,
    block/.style  = {draw, rounded corners=2pt, fill=white, align=center,
                     minimum width=20mm, minimum height=11mm, font=\small},
    flow/.style   = {-{Stealth[length=2.2mm]}, thick},
    group/.style  = {draw=none, rounded corners=4pt, inner sep=3mm},
    glabel/.style = {font=\footnotesize\itshape, anchor=south west},
    dlabel/.style = {font=\scriptsize, midway, above, inner sep=1pt},
]

% ---- Row 1: analogue front end and digital front end -------------------
\node[block] (array) {Antenna\\array};
\node[block, right=of array] (rf)  {RF front end\\(LNA, filter)};
\node[block, right=17mm of rf] (adc) {ADC};
\node[block, right=of adc] (ddc) {Digital\\downconversion};

% ---- Row 2: GPU signal processing and data processing ------------------
\node[block, below=20mm of array] (pc)   {Pulse\\compression};
\node[block, right=of pc]        (dop)  {Doppler\\processing};
\node[block, right=of dop]       (dbf)  {Digital\\beamforming\\(STAP)};
\node[block, right=of dbf]       (cfar) {CFAR\\detection};
\node[block, right=17mm of cfar]  (trk)  {Tracking \&\\data processing};

% ---- Data flow ----------------------------------------------------------
\draw[flow] (array) -- (rf);
\draw[flow] (rf)  -- node[dlabel]{analogue} (adc);
\draw[flow] (adc) -- (ddc);
\draw[flow] (ddc.south) -- ++(0,-6mm) -| node[pos=0.25, below, font=\scriptsize]{per-element I/Q data} (pc.north);
\draw[flow] (pc)  -- (dop);
\draw[flow] (dop) -- (dbf);
\draw[flow] (dbf) -- (cfar);
\draw[flow] (cfar) -- node[dlabel]{detections} (trk);

% ---- Processing domains (drawn behind the blocks) -----------------------
\begin{scope}[on background layer]
  \node[group, fill=gray!15,   fit=(array)(rf)] (gA) {};
  \node[group, fill=blue!10,   fit=(adc)(ddc)]  (gB) {};
  \node[group, fill=green!12,  fit=(pc)(cfar)] (gC) {};
  \node[group, fill=orange!15, fit=(trk)]       (gD) {};
\end{scope}
\node[glabel] at (gA.north west) {Analogue};
\node[glabel] at (gB.north west) {Digital front end (FPGA)};
\node[glabel] at (gC.north west) {CPU / GPU Signal Processing Chain};
\node[glabel] at (gD.north west) {Data processing (CPU)};

% ---- Data-volume annotation ---------------------------------------------
\draw[-{Stealth}, thick, gray]
  ($(gC.south west)+(0,-4mm)$) -- ($(gC.south west -| gD.east)+(0,-4mm)$)
  node[midway, below, font=\footnotesize, text=gray!70!black]
  {data volume decreases along the chain; processing becomes more irregular};

\end{tikzpicture}}
    \caption{Exemplar receive processing chain, with beamforming and subsequent
             signal processing allocated to the GPU. Data volume decreases
             along the chain.}
    \label{fig:processing-chain}
\end{figure}


The data cube is large. A CPI with $N_c=64$ channels,
$N_s=4096$ samples and $N_p=128$ pulses, stored as
single-precision complex values (8\,B per sample),
occupies about 268\,MB. Refreshed every 10\,ms, this is
roughly 27\,GB/s, about 85\% of the 31.5\,GB/s
theoretical peak of a PCIe~4.0~$\times$16 link in each
direction \cite{CITEKEY-AMPERE-WHITEPAPER}. Measured
transfer rates are typically 70--75\% of peak
\cite{venterSOFTWAREDEFINEDPULSEDOPPLERRADAR}, so the host
link alone leaves little headroom. The data volume then
falls along the chain, as beamforming combines the
channels and detection reduces the cube to a list of
detections.


\subsection{GPU Architecture Essentials}
Terminology used follows NVIDIA's CUDA programming model although other vendors provide equivalent concepts. A GPU consists of many streaming mulitprocessors (SMs), each of which executes one or more bvlocks of threads. Threads within a block run in warps of 32 in a single-instruction multiple-threadn (SIMT) model. Rather than minimising the latency of individual threads, GPUs are designed to optimise the overall throughput. It does so by keeping many warps resident and switching between them. The ratio of resident warps to the maximum supported by an SM is known as \emph{occupancy}. High occupancy helps to hide latency but is not a goal in itself: kernels can run faster at lower occupancy when each thread has indepdennt work \cite{volkov}.

In a discrete system, the GPU has its own memory reached over PCIe. Transfers are performed by direct memory access (DMA) from page-locked (pinned) host buffers. It is possible to overlap memory copies with computation by making use of streams. Unified memory is an alternative model where a single address space is presented to the CPU adn GPU and page migration is done on demand, wheras integrated system-on-chip (SoC) share physical DRAM between the two. On the device, a memory hierarchy is present: global memory (HBM or GDDR) is backed by an L2 cache shared by all SMs, and each SM also provides shared memory (controlled by the programmer). Bandwidth falls at each level outwards, most steeply across the host link: on the order of $2TB/s$ for HBM against $32GB/s$ per direction for PCIe4.

Kernls can be either compute-bound or memory-bound. The classfiication that a kernel falls into depends on its operational intensity: the number of floating-point operations it performs per byte of memory traffic. The roofline model \cite{williamsRooflineInsightfulVisual2009} bounds attainable performance by the lower of the peak compute rate and the product of operational intensity and memory bandwidth. The point at which they meet is known as the ridge. For a typical, modern, data-center GPU this sits around $10$FLOP/byte (19.5TFLOPS/s and $\sim$2TB/s for an NVIDIA A100. Kernels below the ridge point are memory bound: making them faster means moving fewer bytes, not performing fewer operations. Section~\ref{sec:workload} applies the roofline model to each stage of hte radar chain.

A GPU is made up of many streaming multiprocessors (SMs). SMs are effectively individual compute units that are responsible for executing a block. Blocks are made up of a group of threads, and are executed in \emph{warps} of 32 threads in lockstep - single-instruction,multiple-thread (SIMT). Rather than minimising per-thread latency, the GPU optimises memory latency by keeping many warps resident and switching between them. The amount of wraps that are resident compared to the maximum number of resident warps is known as \emph{occupancy}.

